% ======================================================================
%  ZANKU / VENTANA SECA — Guía de la solución (para el pitch)
%
%  Documento para el equipo: problemática + evidencia + recorrido
%  comentado de la app (brigadista, jefe, WhatsApp).
%  No sustituye el paper técnico de experimentos.
%
%  Compilar (desde docs/):
%    pdflatex guia-solucion-zanku.tex
%    pdflatex guia-solucion-zanku.tex
% ======================================================================

\documentclass[11pt,a4paper]{article}

\usepackage[utf8]{inputenc}
\usepackage[T1]{fontenc}
\usepackage[spanish,es-tabla]{babel}
\usepackage[a4paper,margin=2.2cm]{geometry}
\usepackage{amsmath,amssymb}
\usepackage{booktabs}
\usepackage{longtable}
\usepackage{array}
\usepackage{tabularx}
\usepackage{xcolor}
\usepackage{enumitem}
\usepackage{graphicx}
\usepackage{titlesec}
\usepackage{fancyhdr}
\usepackage{parskip}
\usepackage{float}
\usepackage{caption}
\usepackage[hidelinks,colorlinks=true,linkcolor=azulvs,urlcolor=azulvs,citecolor=azulvs]{hyperref}

\definecolor{azulvs}{HTML}{1C5FD6}
\definecolor{rojovs}{HTML}{D1382C}
\definecolor{ambarvs}{HTML}{C07A14}
\definecolor{verdevs}{HTML}{177A4A}
\definecolor{grisvs}{HTML}{4A5058}
\definecolor{panelvs}{HTML}{F2F5FA}

\titleformat{\section}{\Large\bfseries\color{azulvs}}{\thesection.}{0.6em}{}
\titleformat{\subsection}{\large\bfseries\color{grisvs}}{\thesubsection}{0.6em}{}
\titleformat{\subsubsection}{\normalsize\bfseries}{\thesubsubsection}{0.6em}{}

\newcolumntype{Y}{>{\raggedright\arraybackslash}X}

\newcommand{\nota}[2]{%
  \par\vspace{4pt}\noindent
  \colorbox{panelvs}{\parbox{\dimexpr\linewidth-2\fboxsep}{%
  \textbf{\color{azulvs}#1.}\ #2}}\par\vspace{4pt}}

\newcommand{\pitch}[1]{%
  \par\vspace{2pt}\noindent
  \textit{\textcolor{azulvs}{Qué decir en el pitch:}\ #1}\par\vspace{4pt}}

\newcommand{\appfig}[3]{%
  \begin{figure}[H]
  \centering
  \includegraphics[width=#2\linewidth]{figuras/app/#1}
  \caption{#3}
  \label{fig:#1}
  \end{figure}}

\graphicspath{{figuras/app/}}

\pagestyle{fancy}
\fancyhf{}
\lhead{\small\textcolor{grisvs}{ZANKU / VENTANA SECA}}
\rhead{\small\textcolor{grisvs}{Guía de pitch}}
\cfoot{\thepage}
\renewcommand{\headrulewidth}{0.4pt}

\begin{document}

% ------------------------------- portada ------------------------------
\begin{titlepage}
\centering
\vspace*{1.5cm}
{\Huge\bfseries\color{azulvs} ZANKU\par}
\vspace{0.4cm}
{\large\textit{App operativa de VENTANA SECA}\par}
\vspace{0.8cm}
{\LARGE\bfseries Guía de la solución\par}
\vspace{0.35cm}
{\large\color{grisvs} Problemática, evidencia y recorrido de la app\\
para armar el pitch del equipo\par}
\vspace{1.2cm}
{\large\textbf{Título de la problemática}\par}
\vspace{0.2cm}
{\large\color{grisvs} «Brigadas antidengue a ciegas: el calendario de\\
cortes de agua que salud no lee»\par}
\vspace{1.5cm}
\begin{tabular}{rl}
\textbf{Desafío:} & Resiliencia climática ante El Niño\\
\textbf{Track:} & 2 --- Salud pública\\
\textbf{Territorio piloto:} & Guayaquil (Guayas), Ecuador\\
\textbf{Producto:} & ZANKU (Next.js + Supabase)\\
\textbf{Proyecto:} & VENTANA SECA\\
\textbf{Fecha:} & 5 de septiembre de 2026\\
\end{tabular}
\vfill
{\small\color{grisvs}
Este documento es el \textbf{libreto} del equipo: explica qué problema
resolvemos, por qué el corte de agua importa, y qué hace cada pantalla
de la app. El paper largo (\texttt{VENTANA-SECA-documento-latex.txt})
guarda la prueba experimental. Toda cifra aquí sale de ese paper, del
código o de fuentes citadas.\par}
\end{titlepage}

\tableofcontents
\newpage

% ======================================================================
\section{Cómo leer este documento}
% ======================================================================

Hay tres capas que no se deben mezclar:

\begin{enumerate}[leftmargin=1.4em]
  \item \textbf{Problemática y evidencia.} Por qué las brigadas llegan
        tarde y por qué el calendario de Interagua es la señal útil.
  \item \textbf{Producto ZANKU.} App de campo (brigadista) + panel de
        mando (jefe) + avisos WhatsApp al vecino.
  \item \textbf{Termómetro nacional.} Pronóstico de casos de dengue
        \emph{del país} (OpenDengue). Corre en el pipeline de
        experimentos; no manda la ruta del barrio.
\end{enumerate}

\nota{En una frase}{Leemos un documento que el Estado ya publica (el
corte de agua) y lo convertimos en la orden de visita de hoy. El dengue
del país se estima aparte. El dengue del barrio se pide al MSP.}

\pitch{«No inventamos un mapa de brotes. Convertimos el comunicado de
corte en la cola de la brigada.»}

% ======================================================================
\section{La problemática}
% ======================================================================

\subsection{Enunciado}

Las brigadas antidengue deciden la ruta con \textbf{casos ya notificados}
y a \textbf{escala gruesa} (cantón / provincia). El calendario de cortes
de agua ---dónde Interagua acaba de pedir \emph{almacenar}--- es
\textbf{público} y \textbf{no entra} a esa programación.

La pregunta operativa no es «¿dónde va a haber dengue?» (no hay $Y$
pública por barrio). Es: \textbf{«¿adónde ir esta semana a tapar
recipientes?»}

\subsection{Quién lo sufre}

\begin{itemize}[leftmargin=1.4em]
  \item El \textbf{jefe de control vectorial}: programa la semana sin
        mirar el comunicado de Interagua.
  \item El \textbf{brigadista}: camina sectores amplios; el tiempo
        escaso no es teclear «3 barriles», es \emph{llegar a la manzana
        correcta}.
  \item El \textbf{vecino}: recibe el aviso de corte («abastézcase»)
        pero no un tip claro de \emph{tapar} el tanque.
\end{itemize}

\subsection{Causa, consecuencia y momento crítico}

\begin{tabularx}{\linewidth}{@{}>{\bfseries}p{3.2cm}Y@{}}
\toprule
Causa & La inspección se decide con información retrasada (casos
notificados) y a escala gruesa.\\
Consecuencia & El recurso llega cuando el criadero ya produjo casos, no
donde van a aparecer criaderos tras el corte.\\
Momento crítico & Los \textbf{7--14 días} posteriores a un corte
$\geq 8$\,h con pedido de almacenar (ciclo huevo$\to$adulto de
\textit{Aedes aegypti}).\\
\bottomrule
\end{tabularx}

\subsection{Cómo se «trata» un corte en Ecuador}

En Guayaquil, Interagua publica el corte (web, X, prensa): horario,
barrios y, a menudo, «abastézcase» / «haga reservas». Si el corte es
largo, mandan tanqueros. Salud no usa ese calendario para programar
brigadas.

En la casa, el protocolo informal es el de siempre: llenar lo que haya.
Quien puede, tanque elevado o cisterna. Quien no, barril, tina o balde.
\textbf{El tanque resuelve el corte; si queda destapado, es el
criadero.}

\nota{Dato del proyecto}{En el CSV de cortes
(\texttt{data/raw/cortes\_interagua\_raw.csv}) la columna
\texttt{pidio\_almacenar} marca cuando el comunicado pide reservas.
Esa bandera dispara la Regla~A.}

\pitch{«Interagua ya dice dónde llenar recipientes. Nosotros convertimos
eso en ``dónde ir a pedir que los tapen''.»}

% ======================================================================
\section{Evidencia: agua, mosquito y dengue}
% ======================================================================

\subsection{Lo que sí está publicado (fuerte)}

\begin{longtable}{@{}p{7.2cm}p{4.8cm}@{}}
\toprule
\textbf{Hallazgo} & \textbf{Fuente}\\
\midrule
\endhead
En Huaquillas (Ecuador), la interrupción del agua fue el \textbf{único}
factor de riesgo significativo para \textit{Ae.~aegypti} en el hogar. &
PLOS NTD (2021)\\
\addlinespace
En Machala, almacenar agua y el acceso irregular predijeron pupas; los
\textbf{barriles} concentraron el 63--69\,\% de las pupas. &
Stewart-Ibarra et al., PLOS ONE (2013)\\
\addlinespace
Riesgo de dengue urbano alto 3--5 meses después de sequía extrema, por
recipientes improvisados; peor donde el desabastecimiento es frecuente. &
Lowe et al., Lancet Planet.\ Health (2021)\\
\bottomrule
\end{longtable}

La cadena que defiende la literatura es:

\begin{center}
\textbf{corte} $\rightarrow$ \textbf{almacenamiento destapado}
$\rightarrow$ \textbf{criadero} $\rightarrow$ \textbf{mosquito}
\end{center}

Por eso ZANKU manda brigada a \emph{tapar recipientes}, no a «predecir
el brote del barrio».

\subsection{Lo que midió este hackatón (honestidad)}

Con la $Y$ pública disponible (casos semanales \textbf{nacionales} /
Guayas frente a cortes de Guayaquil):

\begin{itemize}[leftmargin=1.4em]
  \item Los cortes \textbf{no mejoraron} el pronóstico de casos
        (varios modelos; a veces empeoran el MAE).
  \item 21 rezagos de \texttt{corte\_hn} (0--20 semanas): el IC al
        95\,\% cruza cero en todos.
  \item El rezago de 3--5 meses de Lowe \textbf{no se replicó} a esa
        escala.
\end{itemize}

\nota{Advertencia metodológica}{Huaquillas y Machala miden
\textbf{mosquito en el hogar}. Lowe mide clima$\to$casos a escala de
ciudad. Nosotros medimos \textbf{casos semanales del país} frente a
cortes de una ciudad. Otra $Y$, otra escala: \textbf{no contradice}
esos papers; tampoco los replica.}

\subsection{¿La informalidad valida la propuesta?}

No. Es un \textbf{amplificador}, no la prueba central.

La propuesta no dice «hay informalidad $\Rightarrow$ hay dengue». Dice:
\textbf{corte + pedido de almacenar $\Rightarrow$ recipientes $\Rightarrow$
criadero}. Eso también ocurre en urbanización formal (Urdesa, Alborada):
el comunicado pide reservas y se llena tanque o tina.

Validar de verdad exige $Y$ de barrio: LIRAa / HI--CI--BI
\emph{después} del corte, o dengue por parroquia (MSP). El censo INEC
2022 (\% sin agua entubada, hacinamiento) aún no está integrado; no
inventamos esa cifra.

\pitch{«La evidencia ecuatoriana es agua$\to$mosquito en el hogar. Con
datos públicos no demostramos corte$\to$casos semanales del país. Por
eso el producto es cola operativa, no oráculo de brotes.»}

% ======================================================================
\section{La solución en una frase}
% ======================================================================

\textbf{El comunicado de corte de agua es la orden de visita.}

\begin{figure}[H]
\centering
\footnotesize
\begin{tabular}{@{}c@{}}
Interagua / prensa (corte $\geq 8$\,h + almacenar)\\
$\downarrow$\\
Motor reglas A--D $\rightarrow$ cola priorizada\\
$\downarrow$\\
Jefe: mapa + repartir minizonas\\
$\downarrow$\\
Brigadista: ruta del día + inspección (larvas/pupas)\\
$\downarrow$\\
Panel: HI / CI / BI + cercos · Vecino: WhatsApp\\
\end{tabular}
\caption{Flujo operativo de ZANKU. Paralelo (no dibujado): OpenDengue
$\rightarrow$ termómetro nacional de casos.}
\end{figure}

Dos preguntas, dos $Y$:

\begin{enumerate}[leftmargin=1.4em]
  \item \textbf{Operativa:} ¿adónde ir hoy? $\rightarrow$ reglas sobre
        cortes públicos.
  \item \textbf{Nacional:} ¿cuántos casos tendrá Ecuador la semana que
        viene? $\rightarrow$ GLM-NB / XGBoost sobre OpenDengue.
\end{enumerate}

% ======================================================================
\section{Motor de cola (reglas A--D)}
% ======================================================================

Código: \texttt{web/lib/motor/reglas.ts}. Es un \textbf{proxy operativo},
no un modelo calibrado de dengue por barrio.

\begin{longtable}{@{}cp{3.2cm}Y@{}}
\toprule
\textbf{Regla} & \textbf{Condición} & \textbf{Acción / etiqueta}\\
\midrule
\endhead
\textbf{A} & Corte $\geq 8$\,h, hace 0--14 días, y pidieron almacenar &
\textbf{Cola de esta semana} / «zona riesgosa». Inspeccionar barriles y
tinas; pedir que tapen.\\
\textbf{B} & Corte hace 7--14 días sin cumplir A &
\textbf{Ventana de criadero} / «zona media». Hipótesis biológica;
\textbf{no validada} en Ecuador. Requiere confirmación humana.\\
\textbf{C} & Lluvia / temporada sin corte & Minga de llantas. Sin
catálogo de barrios: \textbf{no genera ítems}.\\
\textbf{D} & Casos cantonales en alza (4 semanas) & Solo \textbf{bonus
de prioridad}. No elige el barrio.\\
\bottomrule
\end{longtable}

Orden de la cola: A $\rightarrow$ B $\rightarrow$ C; dentro de cada
regla, por puntaje (más horas $\Rightarrow$ más arriba; D suma 5 puntos).

\nota{Regla B}{Antes se usó la ventana 12--20 semanas (Lowe). Se cambió
a 7--14 días por el ciclo del vector. Sigue marcada como hipótesis en
código y en la UI.}

\pitch{«La ruta sale del corte $\geq 8$\,h + almacenar, no de un mapa
de probabilidad de brote.»}

% ======================================================================
\section{ZANKU: dos puertas}
% ======================================================================

La portada obliga a elegir rol. No hay un solo login genérico: así el
demo del pitch no mezcla pantallas de campo con las de mando.

\appfig{00-portada.png}{0.42}{Portada de ZANKU: App de campo (brigadista)
y Panel de mando (jefe). Pie: Guayaquil; vecinos con avisos WhatsApp.}

% ======================================================================
\section{Recorrido del brigadista}
% ======================================================================

Menú inferior: \textbf{Mis zonas} · \textbf{Registros} ·
\textbf{Perfil}. La inspección no es pestaña: se abre desde una
minizona.

\subsection{Entrada a campo}

\appfig{01-login-brigadista.png}{0.42}{Login de campo. Demo:
\texttt{brigada.ventana@gmail.com}.}

\subsection{Mis zonas --- la jornada de hoy}

El jefe reparte un \textbf{bloque} de celdas H3 ($\sim$160\,m). El
brigadista solo ve el \textbf{cupo del día: 8 minizonas}
($\sim$40 viviendas). El resto del bloque queda para otros días, para
cubrir el perímetro sin cruzarse con otro compañero.

Qué muestra la pantalla:

\begin{itemize}[leftmargin=1.4em]
  \item Barra de avance (minizonas hechas y viviendas vs meta).
  \item Mapa con paradas numeradas y km de caminata.
  \item Lista: sector, metros a la siguiente, viviendas $n$/meta, si es
        \textbf{cerco} (foco a $\sim$225\,m).
  \item Colores: pendiente, cerco, cubierta.
\end{itemize}

Un toque resalta; dos toques o «Inspeccionar» abre el formulario.

\appfig{02-mis-zonas.png}{0.48}{Mis zonas: cupo diario, mapa de ruta y
lista de paradas (ejemplo Guasmo).}

\pitch{«No mandamos al brigadista a ``todo el sur''. Le damos 8 celdas
de $\sim$160\,m ordenadas para caminar hoy.»}

\subsection{Inspección --- registrar una vivienda}

Formulario en \textbf{4 pasos}. Si la casa está cerrada / renuente /
deshabitada, se guarda solo el estado (cuenta para el denominador tipo
LIRAa) y se salta hogar y recipientes.

\subsubsection{Paso 1 --- Visita}

Estado de la visita, código de vivienda, manzana, GPS. La minizona puede
venir preasignada desde la ruta; el GPS confirma o corrige.

\appfig{03-inspeccion-paso1.png}{0.48}{Paso 1: estado de visita y
contexto del corte (acción + justificación de la Regla~A).}

\subsubsection{Paso 2 --- Hogar}

Habitantes, si tiene red, días sin agua, horas de agua al día, si
almacena, \textbf{motivo} (corte programado / emergente / presión /
costumbre), días con agua almacenada, tanquero.

Aquí se ve la hipótesis en el predio: «¿esta casa llenó recipientes por
el corte?»

\appfig{04-inspeccion-paso2.png}{0.48}{Paso 2: contexto de agua y
almacenamiento en el hogar.}

\subsubsection{Paso 3 --- Recipientes}

Uno por uno: tipo, uso, capacidad, tapado, con agua, ubicación,
\textbf{larvas sí/no}, \textbf{pupas sí/no}, tratamiento. Se pueden
añadir más recipientes.

El brigadista actual ya hace esto a ojo (linterna, cucharón). No hace
falta laboratorio ni clasificador de fotos.

\appfig{05-inspeccion-paso3.png}{0.48}{Paso 3: captura entomológica por
recipiente (tapado, agua, larvas, pupas, tratamiento).}

\subsubsection{Paso 4 --- Acciones}

Checklist de lo ejecutado en el predio: eliminar o tapar recipiente,
aplicar larvicida, instalar malla, entrenar al hogar, entregar
material. Más reinspección. Luego: \textbf{Guardar visita}.

Eso permite comparar, en visitas posteriores, qué paquete de
intervención deja menos criaderos.

Si hay larvas o pupas, el sistema abre un \textbf{cerco} de minizonas
vecinas ($\sim$225\,m) para la ruta.

\appfig{06-inspeccion-paso4.png}{0.48}{Paso 4: checklist de acciones
en el predio y reinspección; botón «Guardar visita».}

\nota{Umbrales OPS}{HI\,4 / CI\,3 / BI\,5. Con pocas casas el número
del panel es ruido; no se vende como LIRAa oficial. La captura
\textbf{sí} construye el dataset que hoy no existe a escala de barrio.}

\subsection{Registros}

Historial de \textbf{sus} visitas, agrupadas por minizona: código,
estado, fecha.

\appfig{07-registros.png}{0.42}{Mis registros: visitas del brigadista
por minizona.}

\subsection{Perfil}

Nombre, rol, brigada. Aviso explícito: \textbf{ZANKU es online} (sin
internet no guarda). Botón Salir.

\appfig{08-perfil.png}{0.42}{Perfil de campo.}

% ======================================================================
\section{Recorrido del jefe de brigada}
% ======================================================================

Menú lateral: \textbf{Resumen} · \textbf{Mapa de riesgo} ·
\textbf{Cola priorizada} · \textbf{Brigada} · \textbf{Alertas}.

\subsection{Entrada a mando}

\appfig{09-login-jefe.png}{0.55}{Login de mando. Demo:
\texttt{jefe.ventana@gmail.com}.}

\subsection{Resumen}

Tablero del día y de las últimas dos semanas:

\begin{itemize}[leftmargin=1.4em]
  \item Cobertura de minizonas; cercos cerrados; focos (viviendas con
        larvas/pupas en 14\,d); brigadistas activos hoy.
  \item Gráfico de viviendas registradas por día (\textbf{esfuerzo de
        campo}, no «bajó el dengue»).
  \item Últimos cortes reportados.
  \item Semáforo \textbf{HI / CI / BI} vs umbrales OPS y \% de positivos
        que son de \textbf{almacenamiento} (prueba directa de la
        hipótesis en campo).
  \item Avance por brigadista y cobertura por sector.
  \item Exportar dataset CSV (nivel recipiente).
\end{itemize}

\appfig{10-panel-resumen.png}{0.95}{Panel Resumen: KPIs, cortes,
índices Stegomyia y avance del equipo.}

\pitch{«El jefe ve esfuerzo y criaderos, no una curva inventada de
casos por manzana.»}

\subsection{Mapa de riesgo}

Mapa de Guayaquil + panel. Círculos de prioridad (zona riesgosa / media)
según el motor. Hexágonos de minizonas y pines de visitas con GPS.

El jefe elige un sector de la cola y puede:

\begin{enumerate}[leftmargin=1.4em]
  \item \textbf{Generar cobertura:} malla de hexágonos; el radio lo pone
        el \textbf{puntaje}, no un clic arbitrario.
  \item \textbf{Repartir entre brigadistas:} bloques contiguos; los
        cercos van primero.
\end{enumerate}

\appfig{11-mapa.png}{0.95}{Mapa de riesgo: sectores priorizados,
minizonas y acciones de generar / repartir.}

\subsection{Cola priorizada}

Lista ordenada que salió de las reglas A/B. Filtros: zona riesgosa, zona
media o todas. Cada ítem: prioridad, nombre, justificación. En B se
avisa que \textbf{aún no entra sola a la ruta de caminata}.

\appfig{12-cola.png}{0.85}{Cola priorizada: sectores A/B con
justificación del motor.}

\subsection{Brigada}

Roster del equipo: teléfono, viviendas totales, si registró hoy, barra
del bloque y chips de la \textbf{ruta de hoy} (cerco en rojo).

\appfig{13-equipo.png}{0.85}{Vista Brigada: avance y ruta del día por
persona.}

\subsection{Alertas --- puente al vecino}

Solo zonas \textbf{A}. Por sector: corte, duración y texto listo.
Botones \textbf{WhatsApp} (abre mensaje) y \textbf{Copiar}.

Ejemplo de mensaje (desde la UI / API): priorizar el barrio, acción de
inspección de recipientes y justificación del corte --- sin prometer
fumigación ni menos dengue.

\appfig{14-alertas.png}{0.85}{Alertas: cortes que generan zona riesgosa
y disparo a WhatsApp.}

% ======================================================================
\section{WhatsApp al ciudadano}
% ======================================================================

El vecino en zona A o B recibe un aviso claro. Los textos
\textbf{no} prometen reducción de dengue ni fumigación confirmada.

\begin{itemize}[leftmargin=1.4em]
  \item \textbf{Zona A} (corte / cola de esta semana): tapa tanques;
        vacía baldes o llantas; pueden pasar brigadistas a
        \textbf{revisar recipientes}.
  \item \textbf{Zona B} (7--14 días post-corte, hipótesis): cuidado en
        casa; B \textbf{no} agenda sola la ruta.
\end{itemize}

Camino técnico: cola A $\rightarrow$ pestaña Alertas o n8n/Twilio
$\rightarrow$ \texttt{GET /api/public/cortes}. Fallback en la app:
\texttt{POST /api/public/whatsapp}.

Texto canónico (código \texttt{web/lib/public/cortes.ts}):

\begin{quote}
\small
\textit{AVISO — Sector (nombre). Se informa que esta semana hay corte
de agua o solicitud de almacenamiento en su sector. Recomendación
sanitaria: cubra tanques y vacíe baldes o llantas. Personal de brigada
podría inspeccionar recipientes. Este comunicado no confirma
fumigación ni reducción de dengue.}
\end{quote}

\appfig{15-whatsapp-mockup.png}{0.55}{Mock de avisos automáticos por
barrio (ZANKU): el mensaje llega; el vecino no tiene que escribir.}

\pitch{«Al vecino le llega el tip el día del corte. A la brigada le
llega la misma zona en la cola.»}

% ======================================================================
\section{Termómetro nacional de dengue}
% ======================================================================

\textbf{Qué hace:} estima cuántos casos de dengue tendrá \textbf{Ecuador}
la semana siguiente, con la serie semanal de OpenDengue / PAHO.

\textbf{Cómo:} train 2021--2022, test 2023.

\begin{itemize}[leftmargin=1.4em]
  \item \textbf{GLM-NB (AR1):} sigue la semana pasada + población.
        MAE $\approx$ 87 casos/semana en 2023.
  \item \textbf{XGBoost:} AR1+AR2+estación+clima. MAE $\approx$ 61.
        Importancia: semana anterior $\sim$45\,\%, estación $\sim$24\,\%.
\end{itemize}

\textbf{Fiabilidad:} bueno como \textbf{carga a 7 días} («el país andará
cerca de $X$ casos»). Flojo como \textbf{aviso de brote} (walk-forward
decae; XGBoost solo marca 1 de 13 picos en ese split). No define barrio.
No usa cortes para decidir la ruta.

\pitch{«Tenemos un termómetro del país donde la $Y$ existe. La ruta del
barrio sale del corte, no de ese modelo.»}

% ======================================================================
\section{Qué sí se defiende / qué no}
% ======================================================================

\begin{longtable}{@{}p{6.2cm}p{6.2cm}@{}}
\toprule
\textbf{Sí (pitch)} & \textbf{No (honestidad)}\\
\midrule
\endhead
Cola operativa Regla~A desde cortes públicos & Predicción de dengue por
barrio / parroquia\\
Captura de recipientes $\rightarrow$ HI/CI/BI & LIRAa OPS con $n$ chico\\
WhatsApp al vecino en zona A/B & Promesa de menos casos o fumigación\\
MAE nacional 2023 (termómetro) & Alerta temprana de brote corroborada\\
Pedido al MSP: dengue o LIRAa por parroquia & Que el rezago Lowe se
replicó aquí\\
App online con roles claros & Funcionamiento offline\\
\bottomrule
\end{longtable}

Límites duros del dato de cortes 2023: de 21 cortes, \textbf{19 no
publicaron horas}; la cola A casi no disparó. El cuello de botella
operativo también es la calidad del comunicado (horas + sectores).

% ======================================================================
\section{Libreto de $\sim$90 segundos}
% ======================================================================

\begin{enumerate}[leftmargin=1.4em]
  \item \textbf{Problema.} Las brigadas van a ciegas respecto del
        calendario de cortes. Salud no lee el aviso que Interagua ya
        publica.
  \item \textbf{Insight.} El corte pide almacenar; el tanque destapado
        es el criadero más productivo y predecible (Huaquillas,
        Machala).
  \item \textbf{Solución.} Regla~A: $\geq 8$\,h + almacenar + 0--14\,d
        $\Rightarrow$ cola de esta semana.
  \item \textbf{Producto.} ZANKU: el brigadista camina 8 minizonas y
        anota larvas; el jefe ve cola, mapa e índices; el vecino recibe
        WhatsApp.
  \item \textbf{Ciencia honesta.} Corte$\to$casos semanales del país: no
        se sostuvo. Corte$\to$visita a tapar recipientes: sí se defiende.
  \item \textbf{Cierre.} «Leemos un documento del Estado y lo convertimos
        en la orden de visita de hoy.»
\end{enumerate}

\vspace{1em}
\noindent\textit{Fin de la guía. En Overleaf: Compiler = pdfLaTeX,
Main document = \texttt{main.tex}. Recompile dos veces para el índice.}

\end{document}
